% !TEX root = ../../../main.tex
% 来源：中学数学实验教材·第二册（几何）
% 章节：实验几何 / 方向、角度与平行
% 原始文件：1.tex，行 1861-1894
% 类型：tikz
% ---
\begin{tikzpicture}
\begin{scope}
\tkzDefPoints{0/1.5/O, -1.5/1.5+1/A, -1.5/1.5-1/B, 1.3/1.5-.5/A', 1.3/1.5+.5/B'}
\tkzDrawSegments(A,O B,O A',O B',O)
\tkzMarkAngles[mark=none, size=.4](A,O,B A',O,B')
\tkzLabelAngle[pos=.6](A,O,B){1}
\tkzLabelAngle[pos=.6](A',O,B'){2}
\node at (0,0){(1)};
\end{scope}
\begin{scope}[xshift=4cm]
\node at (0,0){(2)};
\tkzDefPoints{-1.5/.5/A, 1.5/.5/B, 0/.5/O, -1.2/1.5/A', 1.4/1.5/B'}
\tkzDrawSegments(A,B O,A' O,B')
\tkzMarkAngles[mark=none, size=.5](A',O,A B,O,B')
\tkzLabelAngle[pos=.7](A',O,A){1}
\tkzLabelAngle[pos=.7](B,O,B'){2}
\end{scope}
\begin{scope}[xshift=8cm]
\node at (0,0){(3)};
\tkzDefPoints{-1.5/.4/A, -1.5/1.8/C, 0/1/O, 1.5/.2/D, 1.5/1.6/B, 2/1/E}
\tkzDrawSegments(A,B C,D O,E)
\tkzLabelPoints[left](A,C)
\tkzLabelPoints[right](B,D)
\tkzLabelPoints[below](O)
\tkzMarkAngles[mark=none, size=.45](C,O,A)
\tkzMarkAngles[mark=none, size=.65](D,O,E)
\tkzMarkAngles[mark=none, size=.7](E,O,B)
\tkzMarkAngles[mark=none, size=.3](B,O,C)
\tkzLabelAngle[pos=.6](C,O,A){1}
\tkzLabelAngle[pos=.9](D,O,E){2}
\tkzLabelAngle[pos=.9](E,O,B){3}
\tkzLabelAngle[pos=.5](B,O,C){4}
\end{scope}
\end{tikzpicture}
